\section{完全从零认识一台计算机}

\begin{keyidea}
一台计算机不是一个会自动运行程序的黑盒。它由供电、时钟、处理器、存储器、
互连和外设组成；每一部分只履行自己的有限职责。所谓“启动”，就是这些部分
按照先后依赖逐级建立条件，最终让处理器从一个确定地址读到第一条指令。
\end{keyidea}

这一节不假定读者学过数字电路、计算机组成或硬件史。我们先给每个常用词一个
能操作、能验证的含义，再把它们连成完整机器。后文仍会在具体上下文里重新解释
这些概念；这里的作用不是要求一次记住所有名词，而是让读者知道每个词指向哪一
类真实对象。

\subsection{一、信息怎样变成机器里的 0 和 1}

\topic{bit 是一次二选一}%
\term{bit}{比特、位}是计算机中最小的二值信息单位。一个 bit 只能处于两种
可区分状态之一，通常写成 0 和 1。它可以表示“关/开”“否/是”“未就绪/已就绪”，
也可以只是一个较大数字中的一位。0 和 1 是人给两类物理状态起的名字，不是
漂浮在机器里的抽象小球。

在普通数字电路中，这两类状态主要由\term{电压}{两个位置之间的电势差}范围
表示。例如，某种器件可能把靠近 0 V 的输入认作逻辑 0，把靠近 3.3 V 的输入
认作逻辑 1。实际标准会规定低电平上限、高电平下限以及不可保证的中间区域，
所以“1 就永远等于 3.3 V”并不准确。电路只需可靠区分两个允许范围。

\topic{多个 bit 可以编码更多状态}%
两个 bit 有 \(2^2=4\) 种组合，八个 bit 有 \(2^8=256\) 种组合。八个 bit
通常组成一个\term{byte}{字节}。字节不是“一个字符”的同义词：英文 ASCII
字符常占一个字节，而 UTF-8 编码的汉字通常占三个字节；一条 x86 指令也可能
由一个到十多个字节组成。

\begin{longtable}{L{0.17\linewidth}L{0.24\linewidth}L{0.49\linewidth}}
\toprule
写法 & 含义 & 容易混淆之处 \\
\midrule
\code{b}、bit & 一个二值位 & 网络速率的 Mb/s 通常是每秒多少 megabit \\
\code{B}、byte & 八个 bit & 文件大小的 MB 通常是 megabyte，大小写不能混用 \\
KiB & \(1024\) B & 二进制容量单位；本项目 ROM 的 128 KiB 是 131072 B \\
MiB & \(1024\) KiB & 256 MiB 是 \(256\times1024\times1024\) B \\
GiB & \(1024\) MiB & 本项目容量规格会使用 64 GiB RAM \\
十六进制 & 每一位表示四个 bit & \code{0xFF} 等于二进制 \code{11111111} 和十进制 255 \\
\bottomrule
\end{longtable}

十六进制不是另一种数据，只是更适合查看二进制的写法。一个十六进制数字恰好
覆盖四个 bit，所以两位十六进制数正好描述一个字节。后文看见
\code{0x3F8}、\code{0xFFFF0} 或 \code{0xFFFFFFF0} 时，先把它们理解成
地址或编号的书写形式，不必把它们当作神秘指令。

\subsection{二、CPU 到底是什么}

\term{CPU}{Central Processing Unit，中央处理器}是负责取指令、理解指令并
改变机器状态的电路。它既像计算器，又像交通控制器：算术逻辑单元完成加减、
比较和位运算，控制单元决定下一步读取什么、哪些数据流向哪里。现代 x86 CPU
内部极其复杂，但软件可见的基本契约仍可以从“取指、译码、执行、提交”理解。

\begin{enumerate}[leftmargin=2em]
  \item \textbf{取指}：CPU 用当前指令地址向存储系统请求若干字节。
  \item \textbf{译码}：CPU 按 x86 指令格式判断这些字节表示哪条指令、有哪些操作数。
  \item \textbf{执行}：内部电路完成计算、读写存储器或访问设备。
  \item \textbf{提交}：结果成为软件可见状态，下一条指令地址随之确定。
\end{enumerate}

真实 CPU 会并行、预测和乱序处理许多指令，远比上面四步复杂；但只要没有异常，
它必须让软件观察到与架构规定一致的结果。本书讨论\term{架构}{软件能够依赖
的处理器接口和规则}，而不是尝试复刻某款 Intel 芯片内部的晶体管实现。

\topic{寄存器是 CPU 内部少量而快速的状态格}%
\term{寄存器}{register}直接位于处理器架构状态中。它们数量少、宽度固定，
用于保存当前运算数、地址或控制状态。x86-64 中常见的几类如下。

\begin{longtable}{L{0.20\linewidth}L{0.27\linewidth}L{0.43\linewidth}}
\toprule
寄存器类别 & 例子 & 在完整执行链中的作用 \\
\midrule
通用寄存器 & \code{RAX}、\code{RBX} & 暂存整数、地址、函数参数和计算结果 \\
指令指针 & \code{RIP} & 指向下一段要取出的指令字节；复位时由 CPU 规定初值 \\
栈指针 & \code{RSP} & 指向 RAM 中当前栈顶；调用函数前必须由早期代码建立 \\
标志寄存器 & \code{RFLAGS} & 保存零、进位、中断允许等条件和控制位 \\
控制寄存器 & \code{CR0}、\code{CR3} & 控制执行模式、分页以及当前页表根 \\
段寄存器 & \code{CS}、\code{DS} & 来自早期 x86 分段模型，现代模式仍保留部分职责 \\
\bottomrule
\end{longtable}

寄存器不是 RAM 的另一种名称。CPU 使用寄存器完成当前一步，使用 RAM 保存远
多得多的代码和数据。一次普通加法可能先把 RAM 中的数装入寄存器，完成加法，
再把结果写回 RAM。

\topic{时钟不是“命令 CPU 快一点”的软件开关}%
\term{时钟}{clock}是一种周期性电信号，为同步电路提供状态采样节拍。频率
3 GHz 表示每秒大约有三十亿个周期，但不表示每条指令一定耗一个周期，也不能
直接推出两款 CPU 的实际速度。电源稳定、参考时钟存在、内部锁相环稳定且复位
释放以后，CPU 才能可靠改变状态。

\subsection{三、地址：位置编号不是里面的内容}

\term{地址}{address}是某个可访问位置的编号，类似公寓门牌号。门牌号 203
不是屋内住户，地址 \code{0x7000} 也不是该位置存放的字节。CPU 发出地址、
读写方向和访问宽度，平台再判断哪个部件应该响应。

假设 RAM 中有连续四个字节：

\begin{center}
\begin{tabular}{lllll}
\toprule
地址 & \code{0x7000} & \code{0x7001} & \code{0x7002} & \code{0x7003} \\
\midrule
内容 & \code{0x48} & \code{0x69} & \code{0x21} & \code{0x00} \\
\bottomrule
\end{tabular}
\end{center}

这里每个地址选中一个字节。CPU 从 \code{0x7000} 读取一字节会得到
\code{0x48}；从同一地址读取四字节，会按照 x86 的小端规则把四个相邻字节
组合成一个更大的整数。地址、读取宽度和字节顺序共同决定结果。

\topic{地址空间是一张“哪些编号由谁响应”的地图}%
\term{地址空间}{address space}是所有可表达地址及其规则的集合。一个 64 位
寄存器能写下许多地址，不代表机器真的装有那么多 RAM。平台会把不同范围分给
不同对象：

\begin{itemize}[leftmargin=2em]
  \item 一部分映射到 DRAM，读写会到达内存控制器；
  \item 一部分映射到固件 Flash，读取会返回启动代码；
  \item 一部分映射到设备寄存器，读写会让设备状态机行动；
  \item 一部分没有任何响应者，访问时可能返回固定值、被丢弃或触发故障。
\end{itemize}

因此，“机器有 4 GiB 地址空间”和“机器装了 4 GiB RAM”不是同一句话。设备
窗口和固件窗口也会占用物理地址。进入分页以后，程序使用的
\term{虚拟地址}{先由页表翻译的地址}还要再映射到物理地址；第七章会完整建立
这层关系。

\subsection{四、RAM 是什么：工作时能快速反复读写的主存}

\term{RAM}{Random Access Memory，随机存取存储器}通常指机器运行时保存
代码、变量、栈、页表和缓存的可读写主存。“随机存取”不是指内容随机，而是
强调访问任意位置不必像磁带那样从头顺序走过去；给出地址后，控制器可以直接
访问相应位置。

\topic{为什么 RAM 断电会丢内容}%
个人计算机的主存通常是 \term{DRAM}{Dynamic RAM，动态随机存取存储器}。
每个存储单元利用极小电荷表达状态，电荷会泄漏，所以内存控制器必须周期性
\term{刷新}{重新恢复单元中正在衰减的电荷状态}。断电以后没有刷新和供电，
原有位状态便不再可靠。这种性质叫\term{易失性}{失去供电后内容不能保证保留}。

DRAM 密度高、单位容量成本低，适合做大容量主存，但刚上电不能直接假定可用。
固件通常要配置内存控制器，读取模组参数，并通过
\term{训练}{寻找可靠时序、电压和信号采样参数}让 CPU 与 DRAM 正确通信。
这正是现代实体机器必须先运行一段早期固件，而不能一上电就把复杂内核塞进
普通 RAM 执行的原因之一。

\topic{SRAM 与 DRAM 都是 RAM，但用途不同}%
\term{SRAM}{Static RAM，静态随机存取存储器}只要持续供电便可保持状态，
不需要 DRAM 式刷新，速度也常更高；代价是每一位所需晶体管更多，容量昂贵。
CPU 缓存和芯片内部小存储常用 SRAM，主存则常用 DRAM。这里的“静态”仍不
表示断电保存，它只是说持续供电时不需要周期刷新。

\topic{RAM 里在不同阶段放什么}%
启动以后，同一片 RAM 会被软件划分成不同用途。代码区保存已经装入的机器码；
数据区保存有初值的全局变量；BSS 对应应当清零的未初始化全局变量；栈保存函数
调用现场和局部数据；堆提供动态分配；页表保存地址翻译关系。硬件只看见地址、
字节和访问属性，“这是栈”或“这是页表”是软件建立并维护的含义。

\begin{failurebox}
RAM 可写不代表任何地址都能随便写。地址必须确实映射到已初始化的 RAM，不能
覆盖正在执行的代码、页表或其他所有者的数据。实体机器若 DRAM 尚未训练，
访问主存甚至可能在第一条日志出现之前就失败。
\end{failurebox}

\subsection{五、ROM 是什么：复位时仍能找到的非易失入口}

\term{ROM}{Read-Only Memory，只读存储器}原本指内容在制造时固定、运行时
只能读取的存储器。早期计算机把固定监控程序或启动程序做进掩模 ROM，断电后
内容仍在，CPU 上电时便有确定代码可取。它解决的是一个“先有鸡还是先有蛋”的
问题：要从磁盘读取操作系统需要先运行驱动，而运行驱动本身又需要先有代码。

“ROM”今天常被宽泛地用作“保存启动固件的非易失存储区域”。实体 PC 更常使用
\term{SPI NOR Flash}{通过 SPI 类串行接口访问、能够整块擦除和重新编程的
非易失闪存}。它并非物理上永远只读：在满足解锁、电压和擦写协议时可以更新；
正常运行时，平台会限制写入以避免错误或恶意程序破坏固件。本书继续称其为
ROM，是在描述它在启动地址空间中的角色，不是在断言芯片一定是古老掩模 ROM。

\begin{longtable}{L{0.19\linewidth}L{0.24\linewidth}L{0.25\linewidth}L{0.22\linewidth}}
\toprule
技术 & 内容怎样写入 & 断电后 & 为什么出现 \\
\midrule
掩模 ROM & 芯片制造时固定 & 保留 & 大批量产品成本低，但不能现场升级 \\
PROM & 用户一次性熔写 & 保留 & 允许制造后决定内容，通常不能擦除 \\
EPROM & 紫外线擦除后重写 & 保留 & 开发阶段能够反复修改，但操作不方便 \\
EEPROM & 电方式按字节或小块改写 & 保留 & 在系统中更新少量配置更方便 \\
NOR Flash & 先按块擦除，再编程 & 保留 & 容量、成本与可更新性适合现代固件 \\
\bottomrule
\end{longtable}

\topic{为什么不能只要 ROM，不要 RAM}%
CPU 可以从 ROM 读取并执行指令，却不能把所有运行状态都写回只读区域。函数
调用需要栈，变量会变化，页表和设备队列需要更新，内核也通常远大于最早固件
入口。因此常见链条是：

\[
\begin{aligned}
\text{非易失 Flash 中的固件}
&\rightarrow \text{初始化 RAM}\\
&\rightarrow \text{把后续代码和数据装入 RAM}
\rightarrow \text{跳到 RAM 执行}.
\end{aligned}
\]

本项目为了让复位链可逐字节审计，构建一个 128 KiB 裸 ROM 镜像供 QEMU 映射。
它不是普通应用程序，也不是实体 LattePanda Mu 主板固件的直接替代品。实体
模块内部已有厂商固件和 UEFI 环境；第二章会明确比较两条边界。

\subsection{六、磁盘为什么既不是 RAM，也不是启动 ROM}

\term{磁盘}{这里泛指块存储设备}用于长期保存大容量数据。机械硬盘、SATA SSD、
NVMe SSD、eMMC 和 U 盘的内部技术不同，但操作系统通常通过控制器协议按块请求
数据。它们断电后保留内容，容量大，却不能在 CPU 复位后自动变成可执行指令。

要从磁盘得到一个字节，软件通常必须先知道控制器在哪里，配置命令队列或端口，
发起读取，等待完成，再把返回的数据放进 RAM。因此磁盘中的内核不会“自己跑
起来”。先执行的固件或引导器负责把它读入 RAM、验证格式，然后把控制权交给
内核。

\begin{longtable}{L{0.15\linewidth}L{0.15\linewidth}L{0.18\linewidth}L{0.20\linewidth}L{0.22\linewidth}}
\toprule
对象 & 断电保存 & 运行时可写 & CPU 能否直接按普通地址取指 & 典型职责 \\
\midrule
寄存器 & 否 & 是 & 指令直接使用，不作为主存地址区 & 当前一步的状态 \\
SRAM 缓存 & 否 & 是 & 通常由 CPU 自动管理 & 缓解 CPU 与主存速度差 \\
DRAM 主存 & 否 & 是 & 可以，但须先初始化和映射 & 运行中的代码与数据 \\
固件 Flash & 是 & 受限制 & 平台映射后可以 & 第一段可执行代码 \\
SSD/磁盘 & 是 & 是 & 通常不可以，须经控制器读取 & 内核、程序、文件长期保存 \\
\bottomrule
\end{longtable}

\topic{“内存”这个中文词为什么容易混淆}%
日常语境常把 RAM 叫“内存”，把 SSD 叫“硬盘”。从广义存储层次看，寄存器、
缓存、RAM、ROM 和磁盘都能保存信息；从操作系统语境看，“内存”通常特指
CPU 可寻址的主存 RAM。本书出现歧义时会明确写 DRAM、固件 Flash 或块设备，
而不只说“存储”。

\subsection{七、电压、电流、地和功率：电路图的最低词汇}

\term{电压}{voltage}描述两点之间推动电荷的电势差，单位是伏特 V。说某个
节点是 3.3 V，通常隐含“相对本电路参考地为 3.3 V”。电压必须在器件允许
范围内；把 12 V 直接接到只允许 3.3 V 的输入可能永久损坏器件。

\term{电流}{current}描述单位时间通过截面的电荷量，单位是安培 A。电流必须
沿闭合回路流动：它从电源出发，经负载，再通过地或另一根回流导体回到电源。
只画一根“信号线”而忘记共同参考和回流路径，就没有描述完整的物理通信条件。

\term{电阻}{resistance}描述导体或器件对电流的阻碍，单位是欧姆
\(\Omega\)。在线性近似中，欧姆定律为
\[
V=I R.
\]
\term{功率}{power}描述能量转换速率，单位是瓦特 W；直流近似下
\[
P=V I.
\]
一个 12 V、2 A 的负载约需要 24 W。电源标注“最大 5 A”表示它最多能在规定
条件下提供相应电流，不是它无论连接什么都强行灌入 5 A；实际电流由供电电压、
负载和控制状态共同决定。

\topic{GND 不是把电神秘地吸走的洞}%
\term{GND}{ground，地}在低压数字板上主要是共同电压参考和电流返回网络。
它不一定等同于建筑保护地或机壳地。高速信号的回流会优先沿邻近参考平面流动，
所以原理图上两个端点都叫 GND 只是逻辑连通的声明，PCB 上仍要设计连续、低阻抗
的实际路径。

\topic{电源轨是同一供电网络的名字}%
\code{VDC}、\code{+5V}、\code{+3V3} 都是\term{电源轨}{为一组负载提供
规定电压的网络}。稳压器把输入电压转换为负载需要的电压；保险丝限制严重
过流后继续供能；TVS 管吸收短时过压；肖特基二极管让多个输入汇合时减少反灌。
这些器件不是装饰，每一个都对应正常路径和故障路径。

\begin{failurebox}
本书的实体载板案例用于理解原理图和软件边界，不构成接线安全承诺。真实上电前
必须核对板卡版本、极性、输入电压、电源额定能力和短路情况；示波器地夹、USB
地和外部电源地之间也可能形成危险回路。
\end{failurebox}

\subsection{八、引脚、导线、网络、总线和连接器}

\term{引脚}{pin}是器件对外提供的电气连接点。芯片封装焊盘、插座触点和连接器
金属端子都可以有引脚编号。编号只用于唯一定位，不能凭大小猜功能；功能必须
查器件数据手册和原理图。

\term{网络}{net}是原理图中被声明为电气相连的一组引脚和导线。两段相隔很远
的线只要标有同一个网络名，就表示在 PCB 上必须相连。反过来，图上两条线交叉
但没有连接点，也可能互不相连。\term{导线}{wire}是网络在原理图上的一段画法；
\term{走线}{track/trace}才是 PCB 铜层上的实际几何铜箔。

\term{连接器}{connector}提供可插拔的机械和电气边界。LattePanda Mu 模块
通过 260 个触点的边缘连接器把电源、高速差分信号和低速控制信号交给载板。
“260 针”不代表有 260 种功能：许多触点是地或同一电源，借助并联触点承载
电流、降低阻抗并为高速信号提供返回路径。

\term{总线}{bus}是一组共同遵守某种时序和编码协议的信号及规则。早期并行
总线可能有多根数据线、地址线和控制线；现代 PCIe、USB、SATA 等常用高速
串行差分链路。总线不只是几根铜线，协议还规定谁可以发送、如何表示一位、
怎样确认完成、发生错误如何恢复。

\begin{longtable}{L{0.17\linewidth}L{0.23\linewidth}L{0.50\linewidth}}
\toprule
词语 & 它回答的问题 & 载板案例中的例子 \\
\midrule
引脚 & 器件的哪个接点 & Mu 连接器第 250--260 脚共同提供 \code{VIN} \\
网络 & 哪些接点逻辑相连 & \code{VDC} 把输入汇合点连到模块和多个稳压器 \\
走线 & PCB 上铜怎样铺设 & PCIe 差分对必须按阻抗和长度约束布线 \\
连接器 & 板与板或线缆怎样插接 & HDMI、USB、RJ45、M.2 和 260 针模块座 \\
协议 & 线上电平随时间怎样解释 & I2C 的起始、地址、应答和停止条件 \\
总线/链路 & 哪组线和协议共同传输 & PCIe lane、USB 端口、SPI 固件接口 \\
\bottomrule
\end{longtable}

\subsection{九、单端、差分、时钟和复位信号}

\term{单端信号}{single-ended signal}主要用某根线相对地的电压表达状态。
UART、I2C 和许多控制脚是单端信号。它们仍然需要可靠地参考共同地和回流路径。

\term{差分信号}{differential signal}用一对导线之间的电压差表达信息，通常
标为 P/N、+/- 或 DP/DN。外界噪声若近似同样影响两根线，接收器通过求差可以
抵消一部分干扰。PCIe、USB 3.x、HDMI 和千兆以太网都大量使用差分传输。
差分并不意味着可以随意绕线：两根线要成对、阻抗受控、参考平面连续，长度和
过孔转换也受协议预算约束。

\term{时钟信号}{clock signal}提供周期性时间参考；\term{复位信号}{reset
signal}把状态机带到规定初始状态。复位名后带 \code{\#}、\code{\_N} 或
\code{\_L} 常表示\term{低有效}{拉到逻辑低电平时功能生效}。例如
\code{PLT\_RST\#} 变低意味着平台复位被断言，变高才允许设备离开复位。

\subsection{十、从晶体管到整机：不要混淆层级}

一块现代计算机板可以按下面层级理解：

\[
\begin{aligned}
\text{晶体管}
&\rightarrow \text{逻辑门与存储单元}
\rightarrow \text{功能模块}
\rightarrow \text{芯片}\\
&\rightarrow \text{封装}
\rightarrow \text{计算模块}
\rightarrow \text{载板}
\rightarrow \text{整机}.
\end{aligned}
\]

\term{芯片}{die 或集成电路语境中的 IC}把大量电路集成在半导体上；封装提供
机械保护、散热和可连接焊盘。\term{SoC}{System on Chip，片上系统}把 CPU
核、内存控制器和许多 I/O 控制器集成到一颗芯片。\term{计算模块}{compute
module}再把 SoC、内存、固件存储和必要供电封装成可插拔模块。
\term{载板}{carrier board}为模块提供输入电源转换、连接器、保护器件以及
信号扇出。\term{主板}{motherboard}通常是整机的主要电路板；采用计算模块时，
载板承担了主板的一部分职责。

本书实体案例采用 LattePanda Mu N100/N305 计算模块与 DFR1142 Lite Carrier
V2。CPU 和 DRAM 并不直接画在这张载板原理图上，因为它们位于模块内部；原理图
从 260 针模块连接器开始，完整描述模块与外部电源、HDMI、USB、PCIe、M.2、
以太网、UART、I2C、SPI、RTC 和风扇之间的板级连接。这种“在产品边界上完整”
不同于“把 CPU 芯片内部所有晶体管都展开”。

\subsection{十一、固件、引导器、内核、驱动和应用程序}

这些词都指软件，但处于不同阶段并承担不同权限。

\term{机器码}{machine code}是 CPU 能按架构规则直接译码的指令字节。
汇编语言是机器指令的可读符号写法，C++ 又由编译器转换成机器码。源文件本身
不会被 CPU 直接执行。

\term{固件}{firmware}是与硬件平台紧密结合、通常保存在非易失存储器中的
软件。它负责最早期处理器和内存初始化、发现启动设备并建立向下一阶段交接的
环境。现代 PC 常使用 \term{UEFI}{统一可扩展固件接口}固件；传统 PC 曾广泛
使用 BIOS 接口。本项目则为 QEMU PC 机器模型构建自己的 128 KiB 裸 ROM，
不借用 SeaBIOS 或 OVMF。

\term{引导器}{bootloader}寻找、验证并装入操作系统内核。它可以是固件的一部分，
也可以位于磁盘。本项目 Stage 1 自行用 ATA PIO 从磁盘读取 Kernel 容器，
校验 CRC32 和 ELF64，并准备进入 C++20 内核所需状态。

\term{内核}{kernel}是操作系统中持续以高权限运行的核心，管理 CPU 时间、
内存、异常、中断、设备和进程隔离，并向用户程序提供受控系统调用。
\term{驱动程序}{device driver}把某类设备寄存器和协议封装成内核可使用的
操作。本项目的 UART、PIC/PIT、i8042 和 ATA PIO 代码都是驱动或平台支持代码。

\term{用户程序}{user program}在受限制的 Ring 3 环境运行，不能直接随意
改页表或访问硬件。它通过\term{系统调用}{从用户态受控进入内核的接口}请求
读写文件、创建管道或退出。本项目 Shell 就是 freestanding C++20 用户程序，
不是内核里伪装成命令行的一段函数。

\begin{pdfdiagram}
  \node[os hardware] (power) {电源、时钟\\与复位};
  \node[os hardware,right=of power] (rom) {非易失固件\\ROM/Flash};
  \node[os block,right=of rom] (loader) {引导器\\装入与验证};
  \node[os state,below=of loader] (ram) {RAM\\代码、栈、页表};
  \node[os block,left=of ram] (kernel) {内核\\资源与隔离};
  \node[os state,left=of kernel] (user) {用户程序\\系统调用};
  \draw[os arrow] (power) -- (rom);
  \draw[os arrow] (rom) -- (loader);
  \draw[os arrow] (loader) -- (ram);
  \draw[os arrow] (ram) -- (kernel);
  \draw[os arrow] (kernel) -- (user);
\end{pdfdiagram}
\diagramnote{启动不是“ROM、RAM、磁盘三选一”。固件先从非易失介质执行，初始化
RAM，引导器再把磁盘中的内核装入 RAM，最后由内核管理用户程序。}

\subsection{十二、按下电源键以后，到底发生了什么}

把前面的名词连起来，一次冷启动可以展开为下面这条因果链。不同实体平台的细节
不同，但依赖方向基本一致。

\begin{enumerate}[leftmargin=2em]
  \item 外部电源把能量送入载板。保护、输入选择和稳压电路产生各路电源轨。
  \item 电源监控确认电压进入允许窗口。参考时钟开始振荡并稳定。
  \item 复位电路继续把 CPU 和设备保持在已知状态，直到供电和时钟条件满足。
  \item 复位释放后，CPU 把寄存器设为架构规定值，并从复位向量对应的地址取指。
  \item 平台地址译码把该高端地址送到固件 Flash；Flash 返回机器码字节。
  \item 早期固件建立临时执行环境，初始化内存控制器并完成 DRAM 训练。
  \item 固件枚举必要设备、选择启动项，或把控制权交给自己的下一阶段引导器。
  \item 引导器经存储控制器读取磁盘块，把内核和初始数据复制到 RAM。
  \item 引导器校验文件格式和完整性，建立执行模式、页表、栈及交接信息。
  \item CPU 跳到内核入口。内核接管异常、中断、内存、设备和调度。
  \item 内核装入用户程序，为它建立独立地址空间，通过系统调用提供服务。
\end{enumerate}

其中任何箭头断开，后面都不会自动补救。没有正确电压就没有稳定逻辑；没有时钟
就没有状态推进；没有复位时序就没有可靠初态；没有地址译码就取不到固件；
没有 RAM 初始化就不能建立普通栈；没有存储驱动就读不到内核；没有页表和权限
就不能安全运行用户程序。

\subsection{十三、QEMU 里的“机器”与桌上的实体板}

\term{模拟器}{emulator}用软件实现另一台机器的可见行为。QEMU 的
\code{pc} 机器模型给本项目提供目标 x86-64 CPU、RAM、ROM 映射、ISA 串口、
IDE、PIC、PIT 和 i8042。它不需要真的在宿主桌面上生成 3.3 V 电源轨，也不会
让我们用万用表测到 PCIe 差分对；它直接建立满足架构接口的模型状态。

\begin{longtable}{L{0.19\linewidth}L{0.32\linewidth}L{0.39\linewidth}}
\toprule
观察对象 & 本项目 QEMU PC & LattePanda Mu 实体案例 \\
\midrule
最早固件 & 自研 128 KiB ROM 直接映射 & 模块内部厂商固件，通常提供 UEFI \\
主存 & 命令行指定的虚拟 RAM & 模块上的真实 DRAM，需电源与训练 \\
磁盘接口 & 教学使用传统 IDE/ATA PIO & 常见 NVMe、eMMC、USB 或 SATA 路径 \\
早期输出 & COM1 16550A，宿主标准输入输出 & 需真实 UART 接口或 UEFI 图形输出 \\
中断系统 & PIC 兼容路径和本地 APIC 模型 & 现代 SoC/APIC、ACPI 描述与真实设备拓扑 \\
电气连接 & 由模型代码表达逻辑关系 & 由 4 层 PCB、连接器、保护和电源器件实现 \\
可验证证据 & 串口日志、退出码、ROM 字节、测试 & 电压、波形、枚举信息、固件日志与软件日志 \\
\bottomrule
\end{longtable}

QEMU 不是“虚假的玩具”，实体板也不是 QEMU 的一比一替身。模拟器非常适合把
复位地址、机器码、I/O 协议和失败路径固定下来；实体硬件则暴露电源完整性、
信号完整性、固件兼容性和真实设备枚举。学习时应明确当前问题属于哪一层。

\subsection{十四、硬件历史应该怎样读}

今天的 x86 平台保留许多看似奇怪的机制，原因通常不是设计者故意制造复杂，
而是旧软件兼容具有商业和工程价值。读历史时可以始终追问四件事：

\begin{enumerate}[leftmargin=2em]
  \item 当时要解决什么实际问题？
  \item 当时晶体管数量、内存容量、封装引脚或软件生态有什么限制？
  \item 新机制怎样在这些限制下工作？
  \item 为兼容已有程序，它的哪些痕迹一直保留到今天？
\end{enumerate}

例如，8086 需要用 16 位寄存器表达超过 64 KiB 的地址，于是使用“段值左移四位
再加偏移”的分段形式；后来 80286 和 80386 把段寄存器发展为受保护的描述符
选择子；x86-64 长模式弱化了大部分分段，却仍保留复位、模式切换和
\code{FS}/\code{GS} 等历史接口。后文不会只列年份，而会沿这条“问题
\(\rightarrow\) 约束 \(\rightarrow\) 方案 \(\rightarrow\) 遗留”链解释。

\subsection{十五、遇到陌生名词时的六问法}

后文每遇到 ROM、RAM、寄存器、端口、页表或描述符，都可以用同一组问题拆解：

\begin{enumerate}[leftmargin=2em]
  \item 它是物理器件、线路、数据结构，还是一套协议？
  \item 它由谁创建或初始化，什么时候才有效？
  \item 它保存什么状态，断电或复位后是否还在？
  \item CPU 通过地址、端口还是专用指令访问它？
  \item 正常路径中谁把所有权交给谁？
  \item 条件不满足时会出现什么可观察故障？
\end{enumerate}

以 RAM 为例：它是物理主存及其控制接口；实体平台由固件初始化；保存运行中
代码和数据，断电丢失；CPU 通过物理地址访问；引导器把已装入内核的区域交给
内核；训练失败时可能无法建立栈、无法装入内核或直接复位。能完整回答这六问，
就不再只是背诵“RAM 是随机存取存储器”。

\begin{contractbox}
本书从此处开始采用一个共同写作契约：重要概念不仅给中文和英文全称，还要交代
物理或软件载体、历史动机、建立时机、数据流、所有权与失败表现。读者不需要先
成为硬件专家或历史专家；需要的背景会在第一次依赖它之前建立。
\end{contractbox}
